\clearpage
\section{Scriptural Date and Location}
\zlabel{triptych:concise:chronology:start}
% Generated. Do not edit.
% Deterministic projection of research/chronology.toml.
% schema: 3
% document: liturgy/roman-rite/postconciliar/roman-missal-third-edition-en-us-2011/propers/temporal/pc-s54-twenty-eighth-sunday-in-ordinary-time-year-a
% generated-by: tools/tpt proper-chronology annotations
\newcommand{\chronologyannotationclaim}[6]{#6}
\newcommand{\chronologyannotationcomparisonclaim}[8]{#8}
\newcommand{\chronologyannotationreach}[3]{}
\newcommand{\chronologyannotationgroup}[3]{#3}
\newcommand{\chronologyannotationcomparisongroup}[4]{#4}
\newcommand{\chronologyannotation}[1]{%
  \ifcsname triptychchronologyannotation@#1\endcsname
    \csname triptychchronologyannotation@#1\endcsname
  \else
    \PackageError{triptych}{No chronology annotation for #1}{}%
  \fi}
% appointment-dependency: guidance/liturgy/postconciliar-propers-registry.md sha256=ffe8a6d8b7d79c652bdc49b3d5f6d4dbb32e3506abc7ab0310e745c2b7306ad4
% appointment-dependency: src/gpt/liturgy/roman-rite/postconciliar/roman-missal-third-edition-en-us-2011/propers/registry/README.md sha256=65db141bfd1ec45b556069b256d53a2e65c945e1d7ec9df8ffb07dd76ee0a2bd
% appointment-dependency: src/gpt/liturgy/roman-rite/postconciliar/roman-missal-third-edition-en-us-2011/propers/registry/formula-dispositions.md sha256=c232dc6b4f3fa676c971ba081765471f366de1509dac0ddabca403f012f59879
% appointment-dependency: src/gpt/liturgy/roman-rite/postconciliar/roman-missal-third-edition-en-us-2011/propers/temporal/pc-s54-twenty-eighth-sunday-in-ordinary-time-year-a/instance/manifest.md sha256=684b6a565026a5227b4a1c62817e40a6b9fef727d07d56218ceee78d123422ef
% appointment-dependency: src/gpt/liturgy/roman-rite/postconciliar/roman-missal-third-edition-en-us-2011/propers/temporal/pc-s54-twenty-eighth-sunday-in-ordinary-time-year-a/propers/verified.md sha256=3fc1045aed7de765bb183ff9e3097f7d46f78bb2ba4164e3be84d30042b1208e
% appointment-dependency: src/gpt/liturgy/roman-rite/postconciliar/roman-missal-third-edition-en-us-2011/propers/temporal/pc-s54-twenty-eighth-sunday-in-ordinary-time-year-a/research/chronology-inputs.toml sha256=110e8e186301d2cfe4abece105f7028e7a6d1415445434971af4e7303df4332a
% appointment-dependency: src/gpt/liturgy/roman-rite/postconciliar/roman-missal-third-edition-en-us-2011/propers/temporal/shared/ordinary-time/weeks/28/propers/verified.md sha256=0affca6cbeaee7ee3111b2d8aaa5ad4696fafd00e82673fa32251987095ad657
% appointment-note: entrance: Psalm 129:3-4; identified-basis; vulgate numbering; canonical verse queries, not appointed wording.
% appointment-note: collect: Missal oration without an appointed biblical citation; thematic affinities do not create a Scripture appointment.
% appointment-note: first-reading: Isaiah 25:6-10a; appointed; vulgate numbering; whole-verse query envelope, not subverse chronology or appointed wording.
% appointment-note: psalm: Psalm 23:1-3a,3b-4,5,6; appointed; hebrew numbering; whole-verse query envelope, not subverse chronology or appointed wording.
% appointment-note: psalm: Psalm 23:6cd; response; hebrew numbering; whole-verse query envelope, not subverse chronology or appointed wording.
% appointment-note: second-reading: Philippians 4:12-14,19-20; appointed; vulgate numbering; canonical verse queries, not appointed wording.
% appointment-note: acclamation: Cf. Ephesians 1:17-18; adaptation; vulgate numbering; canonical verse queries, not appointed wording.
% appointment-note: gospel-long: Matthew 22:1-14; appointed; vulgate numbering; canonical verse queries, not appointed wording.
% appointment-note: gospel-short: Matthew 22:1-10; appointed; vulgate numbering; canonical verse queries, not appointed wording.
% appointment-note: offerings: Missal oration without an appointed biblical citation; thematic affinities do not create a Scripture appointment.
% appointment-note: communion-psalm: Cf. Psalm 33:11; adaptation; vulgate numbering; canonical verse queries, not appointed wording.
% appointment-note: communion-john: 1 John 3:2; identified-basis; vulgate numbering; canonical verse queries, not appointed wording.
% appointment-note: after-communion: Missal oration without an appointed biblical citation; thematic affinities do not create a Scripture appointment.
% comparison-dependency: src/gpt/liturgy/roman-rite/postconciliar/roman-missal-third-edition-en-us-2011/propers/temporal/pc-s54-twenty-eighth-sunday-in-ordinary-time-year-a/research/chronology-profile-comparisons.toml sha256=e77bdc39d80aba55b7dbae9c0d2426251d61ca8b0fb7209a94fc65293a23a82b
% entrance: Entrance antiphon
\expandafter\def\csname triptychchronologyannotation@entrance\endcsname{%
  \chronologyannotationgroup{historical-setting}{research-pending}{\textbf{Historical setting} -- Occasion date unresolved.}%
  \space \chronologyannotationgroup{composition}{preferred}{\textbf{Composition}: \chronologyannotationreach{Ps.129.3}{Ps}{true}\chronologyannotationreach{Ps.129.4}{Ps}{true}\chronologyannotationclaim{critical.psalms.latest-composition-boundary}{composition}{catholic-critical-v1}{preferred}{before the Maccabean period, around 165 B.C.; no individual psalm can be dated securely}{Before c. 165 B.C.}}%
}
% first-reading: First reading
\expandafter\def\csname triptychchronologyannotation@first-reading\endcsname{%
  \chronologyannotationgroup{traditional-attribution}{preferred}{\textbf{Traditional attribution}: \chronologyannotationreach{Is.25.10}{Is}{true}\chronologyannotationreach{Is.25.6}{Is}{true}\chronologyannotationreach{Is.25.7}{Is}{true}\chronologyannotationreach{Is.25.8}{Is}{true}\chronologyannotationreach{Is.25.9}{Is}{true}\chronologyannotationclaim{israel.prophets.isaias-traditional-era}{traditional-attribution}{catholic-traditional-v1}{preferred}{from the closing year of Ozias, King of Juda (740 B.C.); no prophecies appear to be later than 701 B.C.}{Isaias (ministry in Souvay's traditional account), B.C. 740--701}.}%
}
% psalm: Responsorial Psalm
\expandafter\def\csname triptychchronologyannotation@psalm\endcsname{%
  \chronologyannotationgroup{historical-setting}{research-pending}{\textbf{Historical setting} -- Occasion date unresolved.}%
  \space \chronologyannotationgroup{composition}{preferred}{\textbf{Composition}: \chronologyannotationreach{Ps.22.1}{Ps}{true}\chronologyannotationreach{Ps.22.2}{Ps}{true}\chronologyannotationreach{Ps.22.3}{Ps}{true}\chronologyannotationreach{Ps.22.4}{Ps}{true}\chronologyannotationreach{Ps.22.5}{Ps}{true}\chronologyannotationreach{Ps.22.6}{Ps}{true}\chronologyannotationclaim{critical.psalms.latest-composition-boundary}{composition}{catholic-critical-v1}{preferred}{before the Maccabean period, around 165 B.C.; no individual psalm can be dated securely}{Before c. 165 B.C.}}%
}
% second-reading: Second reading
\expandafter\def\csname triptychchronologyannotation@second-reading\endcsname{%
  \chronologyannotationgroup{composition}{disputed}{\textbf{Composition} -- disputed: \chronologyannotationreach{Phil.4.12}{Phil}{true}\chronologyannotationreach{Phil.4.13}{Phil}{true}\chronologyannotationreach{Phil.4.14}{Phil}{true}\chronologyannotationreach{Phil.4.19}{Phil}{true}\chronologyannotationreach{Phil.4.20}{Phil}{true}\chronologyannotationclaim{composition.epistle-to-the-philippians}{composition}{catholic-traditional-v1}{disputed}{(Philemon; Colossians; Ephesians; Philippians), 61}{A.D. 61}; \chronologyannotationreach{Phil.4.12}{Phil}{true}\chronologyannotationreach{Phil.4.13}{Phil}{true}\chronologyannotationreach{Phil.4.14}{Phil}{true}\chronologyannotationreach{Phil.4.19}{Phil}{true}\chronologyannotationreach{Phil.4.20}{Phil}{true}\chronologyannotationclaim{composition.epistle-to-the-philippians}{composition}{catholic-traditional-v1}{disputed}{at Rome (A.D. 62-64)}{A.D. 62--64}.}%
}
% acclamation: Gospel acclamation
\expandafter\def\csname triptychchronologyannotation@acclamation\endcsname{%
  \chronologyannotationgroup{composition}{disputed}{\textbf{Composition} -- disputed: \chronologyannotationreach{Eph.1.17}{Eph}{true}\chronologyannotationreach{Eph.1.18}{Eph}{true}\chronologyannotationclaim{composition.epistle-to-the-ephesians}{composition}{catholic-traditional-v1}{disputed}{a period between 58 and 63}{A.D. 58--63}; \chronologyannotationreach{Eph.1.17}{Eph}{true}\chronologyannotationreach{Eph.1.18}{Eph}{true}\chronologyannotationclaim{composition.epistle-to-the-ephesians}{composition}{catholic-traditional-v1}{disputed}{(Philemon; Colossians; Ephesians; Philippians), 61}{A.D. 61}.}%
}
% gospel-long: Gospel: long form
\expandafter\def\csname triptychchronologyannotation@gospel-long\endcsname{%
  \chronologyannotationgroup{narrated-event}{preferred}{\textbf{Event}: \chronologyannotationreach{Matt.22.1}{Matt.22.1-14}{false}\chronologyannotationreach{Matt.22.10}{Matt.22.1-14}{false}\chronologyannotationreach{Matt.22.11}{Matt.22.1-14}{false}\chronologyannotationreach{Matt.22.12}{Matt.22.1-14}{false}\chronologyannotationreach{Matt.22.13}{Matt.22.1-14}{false}\chronologyannotationreach{Matt.22.14}{Matt.22.1-14}{false}\chronologyannotationreach{Matt.22.2}{Matt.22.1-14}{false}\chronologyannotationreach{Matt.22.3}{Matt.22.1-14}{false}\chronologyannotationreach{Matt.22.4}{Matt.22.1-14}{false}\chronologyannotationreach{Matt.22.5}{Matt.22.1-14}{false}\chronologyannotationreach{Matt.22.6}{Matt.22.1-14}{false}\chronologyannotationreach{Matt.22.7}{Matt.22.1-14}{false}\chronologyannotationreach{Matt.22.8}{Matt.22.1-14}{false}\chronologyannotationreach{Matt.22.9}{Matt.22.1-14}{false}\chronologyannotationclaim{life-of-christ.parable-of-the-marriage-feast}{narrated-event}{catholic-traditional-v1}{preferred}{derived A.D. 29 from February, A.U.C. 782- Passover, 782}{The parable of the marriage feast, proposed in the Temple, A.D. 29 (derived)}.}%
  \space \chronologyannotationgroup{composition}{disputed}{\textbf{Composition} -- disputed: \chronologyannotationreach{Matt.22.1}{Matt}{true}\chronologyannotationreach{Matt.22.10}{Matt}{true}\chronologyannotationreach{Matt.22.11}{Matt}{true}\chronologyannotationreach{Matt.22.12}{Matt}{true}\chronologyannotationreach{Matt.22.13}{Matt}{true}\chronologyannotationreach{Matt.22.14}{Matt}{true}\chronologyannotationreach{Matt.22.2}{Matt}{true}\chronologyannotationreach{Matt.22.3}{Matt}{true}\chronologyannotationreach{Matt.22.4}{Matt}{true}\chronologyannotationreach{Matt.22.5}{Matt}{true}\chronologyannotationreach{Matt.22.6}{Matt}{true}\chronologyannotationreach{Matt.22.7}{Matt}{true}\chronologyannotationreach{Matt.22.8}{Matt}{true}\chronologyannotationreach{Matt.22.9}{Matt}{true}\chronologyannotationclaim{composition.gospel-of-matthew}{composition}{catholic-traditional-v1}{disputed}{about A.D. 38-45}{c. A.D. 38--45}; \chronologyannotationreach{Matt.22.1}{Matt}{true}\chronologyannotationreach{Matt.22.10}{Matt}{true}\chronologyannotationreach{Matt.22.11}{Matt}{true}\chronologyannotationreach{Matt.22.12}{Matt}{true}\chronologyannotationreach{Matt.22.13}{Matt}{true}\chronologyannotationreach{Matt.22.14}{Matt}{true}\chronologyannotationreach{Matt.22.2}{Matt}{true}\chronologyannotationreach{Matt.22.3}{Matt}{true}\chronologyannotationreach{Matt.22.4}{Matt}{true}\chronologyannotationreach{Matt.22.5}{Matt}{true}\chronologyannotationreach{Matt.22.6}{Matt}{true}\chronologyannotationreach{Matt.22.7}{Matt}{true}\chronologyannotationreach{Matt.22.8}{Matt}{true}\chronologyannotationreach{Matt.22.9}{Matt}{true}\chronologyannotationclaim{composition.gospel-of-matthew}{composition}{catholic-traditional-v1}{disputed}{about the year 40-42}{c. A.D. 40--42}; \chronologyannotationreach{Matt.22.1}{Matt}{true}\chronologyannotationreach{Matt.22.10}{Matt}{true}\chronologyannotationreach{Matt.22.11}{Matt}{true}\chronologyannotationreach{Matt.22.12}{Matt}{true}\chronologyannotationreach{Matt.22.13}{Matt}{true}\chronologyannotationreach{Matt.22.14}{Matt}{true}\chronologyannotationreach{Matt.22.2}{Matt}{true}\chronologyannotationreach{Matt.22.3}{Matt}{true}\chronologyannotationreach{Matt.22.4}{Matt}{true}\chronologyannotationreach{Matt.22.5}{Matt}{true}\chronologyannotationreach{Matt.22.6}{Matt}{true}\chronologyannotationreach{Matt.22.7}{Matt}{true}\chronologyannotationreach{Matt.22.8}{Matt}{true}\chronologyannotationreach{Matt.22.9}{Matt}{true}\chronologyannotationclaim{composition.gospel-of-matthew}{composition}{catholic-traditional-v1}{disputed}{the years 40-45}{A.D. 40--45}; \chronologyannotationreach{Matt.22.1}{Matt}{true}\chronologyannotationreach{Matt.22.10}{Matt}{true}\chronologyannotationreach{Matt.22.11}{Matt}{true}\chronologyannotationreach{Matt.22.12}{Matt}{true}\chronologyannotationreach{Matt.22.13}{Matt}{true}\chronologyannotationreach{Matt.22.14}{Matt}{true}\chronologyannotationreach{Matt.22.2}{Matt}{true}\chronologyannotationreach{Matt.22.3}{Matt}{true}\chronologyannotationreach{Matt.22.4}{Matt}{true}\chronologyannotationreach{Matt.22.5}{Matt}{true}\chronologyannotationreach{Matt.22.6}{Matt}{true}\chronologyannotationreach{Matt.22.7}{Matt}{true}\chronologyannotationreach{Matt.22.8}{Matt}{true}\chronologyannotationreach{Matt.22.9}{Matt}{true}\chronologyannotationclaim{composition.gospel-of-matthew}{composition}{catholic-traditional-v1}{disputed}{about the year 60-68}{c. A.D. 60--68}; \chronologyannotationreach{Matt.22.1}{Matt}{true}\chronologyannotationreach{Matt.22.10}{Matt}{true}\chronologyannotationreach{Matt.22.11}{Matt}{true}\chronologyannotationreach{Matt.22.12}{Matt}{true}\chronologyannotationreach{Matt.22.13}{Matt}{true}\chronologyannotationreach{Matt.22.14}{Matt}{true}\chronologyannotationreach{Matt.22.2}{Matt}{true}\chronologyannotationreach{Matt.22.3}{Matt}{true}\chronologyannotationreach{Matt.22.4}{Matt}{true}\chronologyannotationreach{Matt.22.5}{Matt}{true}\chronologyannotationreach{Matt.22.6}{Matt}{true}\chronologyannotationreach{Matt.22.7}{Matt}{true}\chronologyannotationreach{Matt.22.8}{Matt}{true}\chronologyannotationreach{Matt.22.9}{Matt}{true}\chronologyannotationclaim{composition.gospel-of-matthew}{composition}{catholic-traditional-v1}{disputed}{about the years 64-67}{c. A.D. 64--67}; \chronologyannotationreach{Matt.22.1}{Matt}{true}\chronologyannotationreach{Matt.22.10}{Matt}{true}\chronologyannotationreach{Matt.22.11}{Matt}{true}\chronologyannotationreach{Matt.22.12}{Matt}{true}\chronologyannotationreach{Matt.22.13}{Matt}{true}\chronologyannotationreach{Matt.22.14}{Matt}{true}\chronologyannotationreach{Matt.22.2}{Matt}{true}\chronologyannotationreach{Matt.22.3}{Matt}{true}\chronologyannotationreach{Matt.22.4}{Matt}{true}\chronologyannotationreach{Matt.22.5}{Matt}{true}\chronologyannotationreach{Matt.22.6}{Matt}{true}\chronologyannotationreach{Matt.22.7}{Matt}{true}\chronologyannotationreach{Matt.22.8}{Matt}{true}\chronologyannotationreach{Matt.22.9}{Matt}{true}\chronologyannotationclaim{composition.gospel-of-matthew}{composition}{catholic-traditional-v1}{disputed}{about the year 50}{c. A.D. 50}.}%
  \space \chronologyannotationcomparisongroup{catholic-critical-v1}{composition}{composition-only}{\textbf{Composition (Catholic critical chronology, for comparison)}: \chronologyannotationreach{Matt.22.1}{Matt}{true}\chronologyannotationreach{Matt.22.10}{Matt}{true}\chronologyannotationreach{Matt.22.11}{Matt}{true}\chronologyannotationreach{Matt.22.12}{Matt}{true}\chronologyannotationreach{Matt.22.13}{Matt}{true}\chronologyannotationreach{Matt.22.14}{Matt}{true}\chronologyannotationreach{Matt.22.2}{Matt}{true}\chronologyannotationreach{Matt.22.3}{Matt}{true}\chronologyannotationreach{Matt.22.4}{Matt}{true}\chronologyannotationreach{Matt.22.5}{Matt}{true}\chronologyannotationreach{Matt.22.6}{Matt}{true}\chronologyannotationreach{Matt.22.7}{Matt}{true}\chronologyannotationreach{Matt.22.8}{Matt}{true}\chronologyannotationreach{Matt.22.9}{Matt}{true}\chronologyannotationcomparisonclaim{critical.gospel-of-matthew}{composition}{catholic-critical-v1}{catholic-critical-v1}{preferred}{passage.united-states-conference-of-catholic-bishops.new-american-bible-revised-edition.english-usccb-web-2026-09-21.matthew-introduction}{post-A.D. 70 date}{The Greek Gospel of Matthew in the NABRE introduction, post-A.D. 70 date}.}%
}
% gospel-short: Gospel: short form
\expandafter\def\csname triptychchronologyannotation@gospel-short\endcsname{%
  \chronologyannotationgroup{narrated-event}{preferred}{\textbf{Event}: \chronologyannotationreach{Matt.22.1}{Matt.22.1-14}{false}\chronologyannotationreach{Matt.22.10}{Matt.22.1-14}{false}\chronologyannotationreach{Matt.22.2}{Matt.22.1-14}{false}\chronologyannotationreach{Matt.22.3}{Matt.22.1-14}{false}\chronologyannotationreach{Matt.22.4}{Matt.22.1-14}{false}\chronologyannotationreach{Matt.22.5}{Matt.22.1-14}{false}\chronologyannotationreach{Matt.22.6}{Matt.22.1-14}{false}\chronologyannotationreach{Matt.22.7}{Matt.22.1-14}{false}\chronologyannotationreach{Matt.22.8}{Matt.22.1-14}{false}\chronologyannotationreach{Matt.22.9}{Matt.22.1-14}{false}\chronologyannotationclaim{life-of-christ.parable-of-the-marriage-feast}{narrated-event}{catholic-traditional-v1}{preferred}{derived A.D. 29 from February, A.U.C. 782- Passover, 782}{The parable of the marriage feast, proposed in the Temple, A.D. 29 (derived)}.}%
  \space \chronologyannotationgroup{composition}{disputed}{\textbf{Composition} -- disputed: \chronologyannotationreach{Matt.22.1}{Matt}{true}\chronologyannotationreach{Matt.22.10}{Matt}{true}\chronologyannotationreach{Matt.22.2}{Matt}{true}\chronologyannotationreach{Matt.22.3}{Matt}{true}\chronologyannotationreach{Matt.22.4}{Matt}{true}\chronologyannotationreach{Matt.22.5}{Matt}{true}\chronologyannotationreach{Matt.22.6}{Matt}{true}\chronologyannotationreach{Matt.22.7}{Matt}{true}\chronologyannotationreach{Matt.22.8}{Matt}{true}\chronologyannotationreach{Matt.22.9}{Matt}{true}\chronologyannotationclaim{composition.gospel-of-matthew}{composition}{catholic-traditional-v1}{disputed}{about A.D. 38-45}{c. A.D. 38--45}; \chronologyannotationreach{Matt.22.1}{Matt}{true}\chronologyannotationreach{Matt.22.10}{Matt}{true}\chronologyannotationreach{Matt.22.2}{Matt}{true}\chronologyannotationreach{Matt.22.3}{Matt}{true}\chronologyannotationreach{Matt.22.4}{Matt}{true}\chronologyannotationreach{Matt.22.5}{Matt}{true}\chronologyannotationreach{Matt.22.6}{Matt}{true}\chronologyannotationreach{Matt.22.7}{Matt}{true}\chronologyannotationreach{Matt.22.8}{Matt}{true}\chronologyannotationreach{Matt.22.9}{Matt}{true}\chronologyannotationclaim{composition.gospel-of-matthew}{composition}{catholic-traditional-v1}{disputed}{about the year 40-42}{c. A.D. 40--42}; \chronologyannotationreach{Matt.22.1}{Matt}{true}\chronologyannotationreach{Matt.22.10}{Matt}{true}\chronologyannotationreach{Matt.22.2}{Matt}{true}\chronologyannotationreach{Matt.22.3}{Matt}{true}\chronologyannotationreach{Matt.22.4}{Matt}{true}\chronologyannotationreach{Matt.22.5}{Matt}{true}\chronologyannotationreach{Matt.22.6}{Matt}{true}\chronologyannotationreach{Matt.22.7}{Matt}{true}\chronologyannotationreach{Matt.22.8}{Matt}{true}\chronologyannotationreach{Matt.22.9}{Matt}{true}\chronologyannotationclaim{composition.gospel-of-matthew}{composition}{catholic-traditional-v1}{disputed}{the years 40-45}{A.D. 40--45}; \chronologyannotationreach{Matt.22.1}{Matt}{true}\chronologyannotationreach{Matt.22.10}{Matt}{true}\chronologyannotationreach{Matt.22.2}{Matt}{true}\chronologyannotationreach{Matt.22.3}{Matt}{true}\chronologyannotationreach{Matt.22.4}{Matt}{true}\chronologyannotationreach{Matt.22.5}{Matt}{true}\chronologyannotationreach{Matt.22.6}{Matt}{true}\chronologyannotationreach{Matt.22.7}{Matt}{true}\chronologyannotationreach{Matt.22.8}{Matt}{true}\chronologyannotationreach{Matt.22.9}{Matt}{true}\chronologyannotationclaim{composition.gospel-of-matthew}{composition}{catholic-traditional-v1}{disputed}{about the year 60-68}{c. A.D. 60--68}; \chronologyannotationreach{Matt.22.1}{Matt}{true}\chronologyannotationreach{Matt.22.10}{Matt}{true}\chronologyannotationreach{Matt.22.2}{Matt}{true}\chronologyannotationreach{Matt.22.3}{Matt}{true}\chronologyannotationreach{Matt.22.4}{Matt}{true}\chronologyannotationreach{Matt.22.5}{Matt}{true}\chronologyannotationreach{Matt.22.6}{Matt}{true}\chronologyannotationreach{Matt.22.7}{Matt}{true}\chronologyannotationreach{Matt.22.8}{Matt}{true}\chronologyannotationreach{Matt.22.9}{Matt}{true}\chronologyannotationclaim{composition.gospel-of-matthew}{composition}{catholic-traditional-v1}{disputed}{about the years 64-67}{c. A.D. 64--67}; \chronologyannotationreach{Matt.22.1}{Matt}{true}\chronologyannotationreach{Matt.22.10}{Matt}{true}\chronologyannotationreach{Matt.22.2}{Matt}{true}\chronologyannotationreach{Matt.22.3}{Matt}{true}\chronologyannotationreach{Matt.22.4}{Matt}{true}\chronologyannotationreach{Matt.22.5}{Matt}{true}\chronologyannotationreach{Matt.22.6}{Matt}{true}\chronologyannotationreach{Matt.22.7}{Matt}{true}\chronologyannotationreach{Matt.22.8}{Matt}{true}\chronologyannotationreach{Matt.22.9}{Matt}{true}\chronologyannotationclaim{composition.gospel-of-matthew}{composition}{catholic-traditional-v1}{disputed}{about the year 50}{c. A.D. 50}.}%
  \space \chronologyannotationcomparisongroup{catholic-critical-v1}{composition}{composition-only}{\textbf{Composition (Catholic critical chronology, for comparison)}: \chronologyannotationreach{Matt.22.1}{Matt}{true}\chronologyannotationreach{Matt.22.10}{Matt}{true}\chronologyannotationreach{Matt.22.2}{Matt}{true}\chronologyannotationreach{Matt.22.3}{Matt}{true}\chronologyannotationreach{Matt.22.4}{Matt}{true}\chronologyannotationreach{Matt.22.5}{Matt}{true}\chronologyannotationreach{Matt.22.6}{Matt}{true}\chronologyannotationreach{Matt.22.7}{Matt}{true}\chronologyannotationreach{Matt.22.8}{Matt}{true}\chronologyannotationreach{Matt.22.9}{Matt}{true}\chronologyannotationcomparisonclaim{critical.gospel-of-matthew}{composition}{catholic-critical-v1}{catholic-critical-v1}{preferred}{passage.united-states-conference-of-catholic-bishops.new-american-bible-revised-edition.english-usccb-web-2026-09-21.matthew-introduction}{post-A.D. 70 date}{The Greek Gospel of Matthew in the NABRE introduction, post-A.D. 70 date}.}%
}
% communion-psalm: Communion: psalm option
\expandafter\def\csname triptychchronologyannotation@communion-psalm\endcsname{%
  \chronologyannotationgroup{traditional-attribution}{disputed}{\textbf{Traditional attribution} -- disputed: \chronologyannotationreach{Ps.33.11}{Ps.33 Ps.39 Ps.70 Ps.85 Ps.94 Ps.97}{true}\chronologyannotationclaim{israel.monarchy.david-traditional-era}{traditional-attribution}{catholic-traditional-v1}{disputed}{reigned from 1055 to 1015 B.C.}{David (reign in the usual chronology), B.C. 1055--1015}.}%
  \space \chronologyannotationgroup{superscription-setting}{preferred}{\textbf{Superscription setting}: \chronologyannotationreach{Ps.33.11}{Ps.33 Ps.55}{true}\chronologyannotationclaim{israel.monarchy.david-at-geth}{superscription-setting}{catholic-traditional-v1}{preferred}{A. M. 2944, A. C. 1060}{David's flight to Achis, king of Geth, A.M. 2944}.}%
  \space \chronologyannotationgroup{composition}{preferred}{\textbf{Composition}: \chronologyannotationreach{Ps.33.11}{Ps}{true}\chronologyannotationclaim{critical.psalms.latest-composition-boundary}{composition}{catholic-critical-v1}{preferred}{before the Maccabean period, around 165 B.C.; no individual psalm can be dated securely}{Before c. 165 B.C.}}%
}
% communion-john: Communion: Johannine option
\expandafter\def\csname triptychchronologyannotation@communion-john\endcsname{%
  \chronologyannotationgroup{composition}{preferred}{\textbf{Composition}: \chronologyannotationreach{1John.3.2}{1John}{true}\chronologyannotationclaim{composition.first-epistle-of-st-john}{composition}{catholic-traditional-v1}{preferred}{from the year 90 to 100 (approximately)}{c. A.D. 90--100}.}%
}

\begin{concisedossiertable}
Psalm & Ps 22:1--6 & Pasture, valley and house: poetic scenes, no located valley. & \chronodate{psalm}{\chronologyannotation{psalm}}\\
\dossierprose{Davidic title, without a dated occasion. The composition bound is the Psalter's broad horizon, not this poem's date (Douay title; NABRE Psalms introduction).}
\midrule
Communion option & Cf. Ps 33:11 & Geth (Gath): superscription setting, not a writing site. & \chronodate{communion-psalm}{\chronologyannotation{communion-psalm}}\\
\cmidrule(lr){2-4}
\dossierevent{Narrated setting: David's flight and assumed madness, before his reign. The title names Achimelech; 1 Kings 21 names Achis. Haydock's A.M. frame dates the narrative, not the psalm.}
\dossierprose{Attribution, setting and composition differ (Douay title; Corbett, ``David, King''; Haydock, 1 Kings chapters 20 through 22; NABRE Psalms introduction).}
\midrule
Entrance & Ps 129:3--4 & The depths signify distress and guilt; no site identified. & \chronodate{entrance}{\chronologyannotation{entrance}}\\
\dossierprose{The Douay text names no author; no particular occasion is established. The composition bound is broad (NABRE Psalms introduction).}
\midrule
First reading & Is 25:6--10a & Zion/Jerusalem: ``this mountain'' follows Is 24:23; all peoples are invited. & \chronodate{first-reading}{\chronologyannotation{first-reading}}\\
\dossierprose{Souvay's traditional ministry era is neither chapter 25's composition date nor the feast's date. Its prophetic referent is an undetermined future (``Isaias,'' Life and First Isaias).}
\midrule
Gospel, long form & Mt 22:1--14 & Temple, Jerusalem: teaching; Antioch: proposed Greek composition setting. & \chronodate{gospel-long}{\chronologyannotation{gospel-long}}\\
Short form & Mt 22:1--10 & Same discourse and composition questions. & \chronodate{gospel-short}{\chronologyannotation{gospel-short}}\\
\cmidrule(lr){2-4}
\dossierevent{Narrated event: teaching after Mt 21:23; Maas's derived date concerns Jesus' Temple discourse, not a historical wedding. Both forms share this dossier.}
\dossierprose{Jacquier finds the early dispersal tradition unreliable, the later reckoning conditional and the Irenaeus inference inconclusive. Durand's early date concerns an Aramaic predecessor; its Greek rendering is undated. Separately, the NABRE introduction favors Greek composition after Jerusalem's destruction, probably at least a decade later, with Antioch plausible. Distinct hypotheses, not a merged range.}
\midrule
Acclamation & Cf. Eph 1:17--18 & Ephesus / wider recipients; writing: Caesarea or Rome proposed. & \chronodate{acclamation}{\chronologyannotation{acclamation}}\\
\dossierprose{Ladeuze and Prat give composition alternatives. Writing place differs from audience; the adaptation has no separate biblical event date.}
\midrule
Second reading & Phil 4:12--14, 19--20 & Philippi: audience; Rome: writing in the traditional accounts. & \chronodate{second-reading}{\chronologyannotation{second-reading}}\\
\dossierprose{Prat and Vander Heeren's Roman-captivity alternatives are historical judgments, not consensus. Hunger and gifts are not independently dated.}
\midrule
Communion option & 1 Jn 3:2 & Writing site unresolved here; future sight has no earthly location. & \chronodate{communion-john}{\chronologyannotation{communion-john}}\\
\dossierprose{The range dates composition, not appearing or vision (Durand, ``The New Testament''; Drum, ``Epistles of Saint John''). Psalm numbers on this sheet are Vulgate.}
\end{concisedossiertable}
\zlabel{triptych:concise:chronology:end}
\clearpage
